% Чуть компактнее вертикальные интервалы, чтобы завершить билет на одной странице.
\begingroup
\setlength{\parskip}{2pt}
\tcbset{before skip=4pt,after skip=4pt,top=5pt,bottom=5pt}
\source{Лекции, с. 39–41; Карпов, §6.4; детали склейки раскрыты составителем}
\begin{definition}
Рёберные компоненты двусвязности — компоненты после удаления всех мостов
(билет 4). Вершинные компоненты двусвязности — блоки (билет 5), включая
мосты с концами и изолированные вершины.
\end{definition}
\begin{lemma}[Склейка в вершине]
Если два планарных графа имеют ровно одну общую вершину и не имеют
общих рёбер, то их объединение планарно.
\end{lemma}
\begin{proof}
Рисунок конечного графа — конечное объединение компактных дуг и точек,
поэтому ограничен. Переносом и гомотетией его можно поместить в любой
заданный открытый диск, сохраняя отсутствие пересечений.
Выберем в каждом рисунке общую вершину \(a\) на внешней грани
(билет 14). Укладку каждого графа можно поместить в замкнутый
топологический диск, граница которого проходит через \(a\) и больше
нигде не касается рисунка. Для этого в секторе внешней грани у \(a\)
проведём два близких непересекающихся выхода к большой окружающей
окружности и соединим их по окружности; полученная простая замкнутая
кривая ограничивает диск с рисунком внутри. Для изолированной вершины
достаточно любого касающегося её диска.
Деформируем диски в два диска, касающиеся только в \(a\), и отождествим
вершины. Внутренности дисков раздельны, новых пересечений нет.
Это также описывает локальную склейку конечного числа рисунков в разных секторах.
\end{proof}
\begin{lemma}[Добавление моста между рисунками]
Два непересекающихся планарных графа можно соединить ребром
между любыми выбранными вершинами с сохранением планарности.
\end{lemma}
\begin{proof}
Делаем обе выбранные вершины внешними, помещаем рисунки в раздельные
диски с этими вершинами на границе и проводим соединяющую дугу
снаружи дисков. Её внутренность не пересекает рисунков.
\end{proof}
\begin{theorem}[Планарность по рёберным компонентам]
Граф планарен тогда и только тогда, когда планарны все его
рёберные компоненты двусвязности.
\end{theorem}
\begin{proof}
Необходимость — наследование планарности подграфами.
После сжатия компонент остаётся лес мостов: цикл в нём дал бы обход
одного из мостов, что противоречит определению моста.
Для каждого дерева этого леса выбираем корень и последовательно
присоединяем рисунки дочерних компонент мостами по предыдущей лемме.
Каждый шаг соединяет два до того раздельных рисунка, поэтому пересечений нет.
Разные деревья леса размещаем отдельно.
\end{proof}
\begin{theorem}[Планарность по блокам]
Граф планарен тогда и только тогда, когда каждый его блок планарен.
\end{theorem}
\begin{proof}
Необходимость следует из планарности подграфов.
Для достаточности сначала рассматриваем связный граф и индукцию
по числу блоков. Если блок один, утверждение дано.
Иначе по дереву блоков (билет 5) существует крайний блок \(B\),
соединённый с остальной частью графа только через одну точку сочленения \(a\).
Остальная часть имеет меньше блоков и планарна по индукции.
Склеиваем её рисунок с рисунком \(B\) в вершине \(a\) по лемме.
Для несвязного графа рисуем его компоненты в непересекающихся дисках.
Изолированные вершины и блоки-мосты заведомо планарны.
\end{proof}
\begin{center}
\begin{tikzpicture}[scale=0.8,every node/.style={circle,fill=cbBlue,inner sep=2pt}]
\node (a) at (0,0) {};\node (b) at (-1,1) {};\node (c) at (-1,-1) {};
\node (d) at (1,1) {};\node (e) at (1,-1) {};
\draw (a)--(b)--(c)--(a)--(d)--(e)--(a);
\node[rectangle,fill=none,above] at (a) {$a$};
\node (f) at (2.7,-1) {};\draw (e)--(f);
\end{tikzpicture}
\end{center}
Здесь два треугольника — разные вершинные блоки, но лежат в одной
рёберной компоненте двусвязности. Последнее ребро — мост.

\endgroup
